\section{Stopping: from a Crossing to a Target}
\label{sec:stopping}

The law \eqref{eq:cert} estimates the error at a given $K$; a practitioner needs a
decision, and there are two ways to read it as a rule.  The first asks where the precision
term overtakes the exact term,
\begin{equation}
K_{\mathrm{cross}}(u) \;=\; \min\{K:\ E_K \le u\,S_{\mathrm{pert}}(K)\},
\label{eq:cross}
\end{equation}
the reading the two-term structure invites; as a stopping time it is also wrong by up to
three orders of magnitude in error (Table~\ref{tab:cross}, appendix).  The second asks which $K$ first meets a stated accuracy,
\begin{equation}
K(\varepsilon) \;=\; \min\{K:\ J(K) \le \varepsilon\},
\qquad J(K) = E_K + c\,u\,S_{\mathrm{pert}}(K),
\label{eq:target}
\end{equation}
and it is the one we recommend and verify.

\begin{theorem}[Stopping at a target accuracy]
\label{thm:stop}
Let $J(K) = E_K + c\,u\,S_{\mathrm{pert}}(K)$ with the calibrated $c$ and $K_{\max}$ the
largest step count the caller will run.  If $\varepsilon > \min_{K \le K_{\max}} J(K)$
then $K(\varepsilon)$ exists and is finite; if $\varepsilon$ is below the predicted floor
the rule returns nothing, and no step count exists.  At the returned step count
$\mathrm{err}(K(\varepsilon)) \le 1.16\,\varepsilon$ in all 62 cells tested (median
$0.45\,\varepsilon$), and it is below $\varepsilon$ in 60 of 62.
\end{theorem}

Existence is immediate; the content is the last sentence, measured over
$d\in\{64,128\}$, three $\ell$, four $p$ and eight
targets from $0.3$ to $10^{-4}$ (Table~\ref{tab:stop}); at $\ell=10^{-3}$, $p=10$ the
predicted floor, $0.0845$, takes the refusal branch.

\begin{table}[t]
\centering
\caption{Target-accuracy stopping rule \eqref{eq:target} against measurement.  Eight
representative rows of the 62 defined cells.  Ratio $=$ measured error / target; the
refusal branch is taken exactly when the target is below the predicted floor
$\min_K J(K)$.}
\label{tab:stop}
\begin{tabular}{rrrr rc}
\toprule
$\ell$ & $p$ & $\varepsilon$ & $K(\varepsilon)$ & measured & ratio \\
\midrule
$10^{-2}$ & 12 & $0.3$  & 3 & \e{2.80}{-1} & 0.93 \\
$10^{-2}$ & 12 & $0.1$  & 4 & \e{1.67}{-2} & 0.17 \\
$10^{-2}$ & 12 & $0.01$ & 5 & \e{3.46}{-3} & 0.35 \\
$10^{-3}$ & 12 & $0.3$  & 5 & \e{1.36}{-1} & 0.45 \\
$10^{-3}$ & 12 & $0.1$  & 6 & \e{1.89}{-2} & 0.19 \\
$10^{-3}$ & 16 & $0.003$& 6 & \e{2.21}{-3} & 0.74 \\
$10^{-4}$ & 16 & $0.3$  & 7 & \e{4.50}{-2} & 0.15 \\
$10^{-4}$ & 16 & $0.03$ & 8 & \e{9.55}{-3} & 0.32 \\
\midrule
\multicolumn{3}{r}{summary, 62 cells} & \multicolumn{2}{c}{ratio $= 0.020$--$1.153$, median $0.451$} \\
\multicolumn{3}{r}{within $1.5\times$} & \multicolumn{2}{c}{62/62} \\
\multicolumn{3}{r}{below the target} & \multicolumn{2}{c}{60/62} \\
\bottomrule
\end{tabular}
\end{table}

The rule is not monotone in $\varepsilon$: over the 44 consecutive
target pairs $K(\varepsilon)$ moves in the expected direction 26 times, because $J$ is not
monotone in $K$; both members of an inverted pair still meet their own targets, and a
$3.5\times$ over-estimate of $c$ changes no selected $K$.

\paragraph{Why the crossing rule is the wrong stopping time.}
Equation~\eqref{eq:cross} does describe where the floor takes over, at the $u$-scale
\eqref{eq:cert} predicts, but it does not locate the argmin: on the frozen real momentum
spectra ($\ell=10^{-3}$, $K$ swept to 10) it stops one to three steps early --- at $p=8$
where the error is $9.1\times$ its value at $K=6$, and $926\times$ in the wide window at
$p=24$ --- because it asks where the two terms are equal in size, where the sum is still
dominated by the falling term, and the error is minimized where they are equal in
\emph{slope}; stopping late is cheap ($6.237\times10^{-2}$ at $K=10$ against
$6.229\times10^{-2}$ at $K=6$).

\paragraph{A lower bound on the dynamic range.}
\label{sec:ellsens}
The last input is $\ell$, which a caller does not have.  With the design and the rule
given $\hat\ell=\alpha\ell$, over 36 cells ($d\in\{64,128\}$,
three $\ell$, $p\in\{8,12\}$) at three targets each
($\varepsilon=r\min_K J(K)$, $r\in\{1.2,3,10\}$; Table~\ref{tab:ell}),
under-estimating is safe and costs arithmetic: at $\alpha\le0.5$ no target is missed, the
error stays below $0.99\,\varepsilon$, and the rule either returns a step count within
three steps of the right one or refuses a target it could have met (8 to 16 of 36 cells,
all at $r=1.2$).  Over-estimating is neither safe nor visible: at $\alpha=10$, 24 targets
are missed with errors up to $244\,\varepsilon$ and the caller's own prediction
under-states the measurement by up to $1275\times$ ($\e{6.1}{-4}$ against $0.782$ at
$\ell=10^{-2}$, $p=12$).  A design fitted to $[\hat\ell,1]$ leaves the polynomial
outside its minimax interval below $\hat\ell$, so the orbit converges to the wrong point.
So $\sigma_{\min}$ should be passed as a lower bound --- a power iteration count with a
safety factor.
